\documentclass[11pt]{article}
\usepackage[margin=1in]{geometry}
\usepackage{amsmath, amssymb}

\title{Exercise 3.7: Maze Robot Reward Design}
\author{Sutton \& Barto, Section 3.2}
\date{}

\begin{document}
\maketitle

\section*{Problem}

Imagine that you are designing a robot to run a maze. You decide to give it a
reward of $+1$ for escaping from the maze and a reward of zero at all other times.
The task seems to break down naturally into episodes---the successive runs through
the maze---so you decide to treat it as an episodic task, where the goal is to
maximize expected total reward~(3.7). After running the learning agent for a while,
you find that it is showing no improvement in escaping from the maze. What is going
wrong? Have you effectively communicated to the agent what you want it to achieve?

\section*{Solution}

\subsection*{What's going wrong}

With undiscounted episodic returns (equation~3.7), the return from any trajectory
that eventually reaches the exit is:
\[
    G = 0 + 0 + \cdots + 0 + 1 = 1
\]
regardless of how many steps it takes. A path that escapes in 10 steps has the
same return as one that wanders for 10{,}000 steps. The agent has no incentive to
escape \emph{quickly}---only to escape \emph{eventually}. Since even a random
policy will eventually escape (assuming the maze is solvable), all policies look
equally good.

\subsection*{The core issue}

The reward signal does not communicate \textbf{urgency}. The agent's goal is to
maximize expected return, but every successful trajectory yields the same return
of~$+1$. There is no gradient for the agent to follow toward faster escape.

\subsection*{Fixes}

Any of these would work:

\begin{enumerate}
    \item \textbf{Use discounting ($\gamma < 1$)}: The return becomes
          $G = \gamma^{K-1}$, where $K$ is the number of steps to escape.
          Faster escape $\to$ higher return.
    \item \textbf{Give $-1$ reward every step} (instead of $+1$ at the end):
          The return becomes $G = -K$, so the agent is directly penalized for
          every extra step. This is what the textbook recommends for maze tasks.
\end{enumerate}

Either fix is sufficient. Note that combining $-1$ per step with discounting
would also work, but discounting actually \emph{weakens} the penalty for long
episodes: the return $G = -(1 + \gamma + \cdots + \gamma^{K-1}) = -(1-\gamma^K)/(1-\gamma)$
is bounded by $-1/(1-\gamma)$, whereas the undiscounted $G = -K$ grows without bound.
So for this episodic task, the simplest and most effective fix is option~2: reward
$= -1$ per step, no discounting.

All approaches communicate the same thing: \emph{get out as fast as possible}.
The original $+1$-at-exit reward only says ``get out'' without the
``as fast as possible'' part.

\end{document}
